
\section{Related Work}

\subsection{Novel View Synthesis}

Neural Radiance Fields~\cite{mildenhall2020nerf,barron2021mip} achieve photo-realistic rendering quality by applying volume rendering which aggregates the sampled neural features along the ray. 
%
Recently, 3D Gaussian Splatting~\cite{kerbl2023_3dgs} has made significant advances in neural rendering for its real-time rendering efficiency.
%
This outstanding technique models static scenes by optimizing a set of Gaussian primitives~\cite{kerbl2023_3dgs, lu2024scaffold, yu2024mip}, including properties of position, scaling, rotation, and opacity. 
%
Some methods progress further to model dynamic scenes with time-varying Gaussian primitives~\cite{li2024spacetime, yang2024deformable, yan2024street, sun2025splatter}, with 4D neural representation~\cite{Wu2024_4dgs, xu20244k4d, xu2024representing}, or with on-the-fly streamable training~\cite{luiten2023dynamic, sun20243dgstream, girish2025queen, gao2025hicom}.
%
Although immensely accelerated, the inevitable per-scene optimization still requires minutes to accomplish.

To eliminate the long-time optimization process, generalizable neural rendering methods~\cite{wang2021ibrnet, yu2021pixelnerf, charatan2023pixelsplat, zheng2024gpsgaussian} have been developed for feed-forward novel view rendering.
%
Typically, these generalizable paradigms resort to leveraging the learned 2D/3D priors from extensive data~\cite{zhou2018stereo, liu2021infinite, yu2021function4d} to ease the long-term optimization.
%
In this practice, ENeRF~\cite{lin2022enerf} integrates cost volume to provide a coarse depth initialization and thus reduce the sampling points, leading to an efficient framework.
%
With respect to Gaussian-based methods~\cite{xu2024depthsplat, charatan2023pixelsplat, zheng2024gpsgaussian, liu2024mvsgaussian}, the solution is parallel to that in generalizable NeRF, \textit{e.g.} by using epipolar stereo and cost volume.
%
However, some methods~\cite{zheng2024gpsgaussian} necessitate ground truth depth for training, and others~\cite{charatan2023pixelsplat, chen2024mvsplat, zhou2024gps} are limited under sparse input views, due to the difficulty of establishing correlation with small overlap of input views.


\subsection{3D Reconstruction}

Multi-view stereo is a traditional 3D reconstruction technique, which can be categorized according to the output modality, including point cloud~\cite{lhuillier2005quasi, furukawa2009accurate}, volumetric representation~\cite{seitz1999photorealistic, kutulakos2000theory}, and depth map~\cite{campbell2008using, schonberger2016pixelwise}.
%
MVSNet~\cite{yao2018mvsnet} opens up the era of deep learning-based MVS methods.
%
Binocular Stereo~\cite{zabih1994non}, as a special kind of MVS, aims to find the maximum correspondence on the horizontal epipolar line.
However, such methods struggle with invisible issues and large sparsity of input views.
%
Progressing further, the neural implicit surface methods~\cite{wang2021neus, wang2023neus2, li2023neuralangelo}, as a variant of neural radiance fields, perform accurate 3D reconstruction with only rendering loss, which avoids collecting 3D properties for training.
%
With the prevalence of 3D Gaussian Splatting in NVS, a body of research~\cite{huang20242d, dai2024high, guedon2024sugar, lyu20243dgsr} attempts to adapt it to multi-view 3D Reconstruction with flattened or surfel shaped Gaussian primitives.
%
These neural rendering based methods typically rely on dense input views to supervise with rendering loss and long-time optimization.


More recently, DUSt3R\cite{wang2024dust3r} proposes a novel 3D representation, defining point maps on a pair of source views, aligning a pixel to a free 3D point, bypassing the need for camera poses.
%
In such ill-posed conditions, both DUSt3R and its follow-ups~\cite{leroy2024grounding, wang2024moge} bound point maps in scale-invariant canonical space.
%
NoPoSplat~\cite{ye2024nopose} and Splat3R~\cite{smart2024splatt3r} incorporate such representation into the rendering pipeline for static scenes.
%
Some followers~\cite{lu2024align3r, zhang2024monst3r} point out that when handling dynamic scenes, DUSt3R encounters two limitations: (1) the misaligned background points, and (2) incorrect foreground depth estimation, causing some regions placed in the background.
%
They address these issues with a global test time optimization on the whole video, while we commit to probing a feed-forward solution in this paper.
